\subsubsection{Graph examples}\label{methods:graph_examples}

The following examples instantiate the formal patterns with observations
from the resource. Table~\ref{tab:prefixes} defines the namespaces used
in the Turtle examples and SPARQL queries. Prepend these declarations
when parsing an example. The listing validator compares each example
with the corresponding modules and pinned reference terms.

\begin{table}[h]
\centering
\caption{Prefix declarations used in all Turtle and SPARQL listings.}
\label{tab:prefixes}
\small
\begin{tabular}{ll}
\toprule
\textbf{Prefix} & \textbf{Namespace} \\
\midrule
\texttt{rak:}      & \texttt{https://rubalkhali.science/kb/RAK\_} \\
\texttt{rak\_kb:}  & \texttt{https://rubalkhali.science/kb/} \\
\texttt{sio:}      & \texttt{http://semanticscience.org/resource/SIO\_} \\
\texttt{owl:}      & \texttt{http://www.w3.org/2002/07/owl\#} \\
\texttt{pato:}     & \texttt{http://purl.obolibrary.org/obo/PATO\_} \\
\texttt{uo:}       & \texttt{http://purl.obolibrary.org/obo/UO\_} \\
\texttt{envo:}     & \texttt{http://purl.obolibrary.org/obo/ENVO\_} \\
\texttt{ncbitaxon:}& \texttt{http://purl.obolibrary.org/obo/NCBITaxon\_} \\
\texttt{obo:}      & \texttt{http://purl.obolibrary.org/obo/} \\
\texttt{rdf:}      & \texttt{http://www.w3.org/1999/02/22-rdf-syntax-ns\#} \\
\texttt{rdfs:}     & \texttt{http://www.w3.org/2000/01/rdf-schema\#} \\
\texttt{xsd:}      & \texttt{http://www.w3.org/2001/XMLSchema\#} \\
\texttt{geo:}      & \texttt{http://www.opengis.net/ont/geosparql\#} \\
\texttt{dcterms:}  & \texttt{http://purl.org/dc/terms/} \\
\bottomrule
\end{tabular}
\end{table}

Listing~\ref{lst:ttl_sampling} represents the Trip~1, Site~1
shallow-subsurface specimen. The sampling process targets a site and
outputs a collection containing the specimen. The specimen also links
directly to its source site, allowing queries to recover provenance
without traversing the sampling event.

\begin{lstlisting}[caption={Turtle representation of a sampling process.}, label={lst:ttl_sampling}, basicstyle=\ttfamily\scriptsize, frame=single]
rak:P100001 a rak:0000024 ; # Sampling Process
  sio:000291 rak:1000001 ; # hasTarget
  sio:000229 rak:7000001 . # hasOutput (Collection)

rak:7000001 a sio:001419 ; # Material sample collection
  sio:000059 rak:6000001 . # hasMember

rak:6000001 a rak:0000021 ; # Deep Soil Sample
  sio:000244 rak:1000001 . # isDerivedFrom
\end{lstlisting}

The same specimen is the input of the extraction process in
Listing~\ref{lst:ttl_protocol}. A protocol specification identifies the
procedure, and the output DNA extract carries a concentration quality.
A separate measurement value quantifies that quality through the
reciprocal SIO measurement relations.

\begin{lstlisting}[caption={Turtle representation of a DNA extraction protocol.}, label={lst:ttl_protocol}, basicstyle=\ttfamily\scriptsize, frame=single]
rak:P200001 a sio:000994 ; # Scientific Protocol Execution
  sio:000339 rak_kb:DNA_ext_powersoil ; # isSpecifiedBy
  sio:000230 rak:6000001 ; # hasInput (Soil Sample)
  sio:000229 rak:6800001 . # hasOutput (DNA Extract)

rak:6800001 a rak:0000040 ; # DNA Extract
  sio:000008 rak:5900001 . # hasAttribute (Concentration Quality)
\end{lstlisting}

Taxonomic observations bind a count to both a taxon and its source
sequence dataset (Listing~\ref{lst:ttl_tax}). Here, run
\texttt{ERR16783308}, profile \texttt{V10Dr1}, has 853,814 reads
assigned to Bacteria. Each taxonomic dataset contains one rank and
one abundance basis; count and relative-abundance datasets retain the
same classifier lineage and source-profile context.

\begin{lstlisting}[caption={Turtle representation of taxonomic abundance formalized from a feature table.}, label={lst:ttl_tax}, basicstyle=\ttfamily\scriptsize, frame=single]
rak:P290001 a rak:0000071 ; # Bioinformatic Workflow
  sio:000339 rak:L000012 ; # 16S amplicon processing protocol
  sio:000230 rak:7791009 ; # hasInput (FASTQ Dataset, ERR16783308)
  sio:000229 rak:7740001 . # hasOutput (Taxonomic Dataset)

rak:7740001 a rak:0000074 ; # taxon absolute abundance dataset
  sio:000059 rak:V00000001 . # hasMember (Observation Value)

rak:V00000001 a rak:0000076 ; # Absolute abundance value
  sio:000215 rak:Q00000001 ; # isMeasurementValueOf
  rak:2000025 "Domain: Bacteria" ; # classifier lineage
  rak:2000026 "853814.0"^^xsd:double . # absolute read count

rak:Q00000001 a rak:0000078, pato:0000070 ; # Absolute Abundance Quality
  sio:000011 ncbitaxon:2, rak:7791009 ; # taxon and FASTQ bearers
  sio:000216 rak:V00000001 . # hasMeasurementValue
\end{lstlisting}

Laboratory XRF specializes the measurement pattern by linking the
process to its archived specimen. Listing~\ref{lst:ttl_xrf} shows the
aluminium observation for \texttt{V1Sr1}. Its numerical value and
quality have reciprocal links. The value retains the instrument-reported
magnitude with undocumented concentration units. Field sessions instead
identify a target site through \texttt{sio:000291}.

\begin{lstlisting}[caption={Turtle representation of an XRF geochemistry measurement.}, label={lst:ttl_xrf}, basicstyle=\ttfamily\scriptsize, frame=single]
rak:P300072 a rak:0000025 ; # XRF Measuring Process
  sio:000230 rak:6001732 ; # hasInput (sample V1Sr1)
  sio:000339 rak:L_Surface_Lab ; # laboratory protocol
  sio:000229 rak:4302121 . # hasOutput

rak:4302121 a rak:0000507 ; # Al concentration value
  sio:000215 rak:5302121 ; # isMeasurementValueOf
  rak:2000012 "0.78"^^xsd:double . # Concentration

rak:5302121 a rak:0000107 ; # Al concentration quality
  sio:000011 rak:6001732 ; # isAttributeOf
  sio:000216 rak:4302121 . # hasMeasurementValue
\end{lstlisting}

Shared inputs and outputs connect extraction, library preparation and
sequencing (Listing~\ref{lst:ttl_sequencing}). The final sequence
dataset links to its ENA accession. A SPARQL traversal can therefore
recover the recorded provenance from specimen to sequence run.

\begin{lstlisting}[caption={Turtle representation of the sequencing workflow from soil sample T1Dr1 to ENA run ERR16061083.}, label={lst:ttl_sequencing}, basicstyle=\ttfamily\scriptsize, frame=single]
rak:P200923 a sio:000994 ; # DNA Extraction Process
  sio:000339 rak_kb:DNA_ext_powersoil ; # isSpecifiedBy
  sio:000230 rak:6000577 ; # hasInput (T1Dr1 Soil Sample)
  sio:000229 rak:6800540 . # hasOutput (DNA Extract)

rak:P250001 a rak:0000065 ; # Library Preparation Process
  sio:000339 rak:L000010 ; # Illumina 16S Prep Guide
  sio:000230 rak:6800540 ; # hasInput (DNA Extract)
  sio:000229 rak:6690001 . # hasOutput (Amplicon Library)

rak:P280001 a rak:0000066 ; # Sequencing Process
  sio:000230 rak:6690001 ; # hasInput (Amplicon Library)
  sio:000229 rak:7790001 . # hasOutput (FASTQ Dataset)

rak:7790001 a rak:0000063 ; # FASTQ Dataset
  rdfs:seeAlso <https://identifiers.org/insdc.run/ERR16061083> .
\end{lstlisting}

Quality-control observations retain the read direction through separate
forward and reverse metric classes. Listing~\ref{lst:ttl_qc} gives the
forward-read GC content for run \texttt{ERR16062319}: 58.0\%.
The quality belongs to the same FASTQ dataset consumed by the
quality-control process.

\begin{lstlisting}[caption={Turtle representation of forward-read GC content for ENA run ERR16062319.}, label={lst:ttl_qc}, basicstyle=\ttfamily\scriptsize, frame=single]
rak:P350001 a rak:0000200 ; # Sequencing QC Process
  sio:000230 rak:7790263 ; # hasInput (FASTQ Dataset)
  sio:000229 rak:4450001 . # hasOutput (QC Value)

rak:4450001 a rak:0000217 ; # Forward GC Content Value
  sio:000215 rak:5550001 ; # isMeasurementValueOf
  sio:000221 uo:0000187 ; # Percent
  rak:2000074 "58.0"^^xsd:double . # has GC content

rak:5550001 a rak:0000205 ; # Forward GC Content Quality
  sio:000011 rak:7790263 ; # isAttributeOf (FASTQ Dataset)
  sio:000216 rak:4450001 . # hasMeasurementValue
\end{lstlisting}
